In this paper, I explore the long-term migration patterns of Japanese Americans
who were forcibly relocated during World War II.
Specifically, I investigate whether geographic proximity to the internment camps
established by the War Relocation Authority (WRA)
influenced migration rates among the Japanese American population
in the decades following internment.
The central research question is
whether counties closer to these internment camps experienced a disproportionately higher rate of Japanese American migration compared to more distant counties.

The US government's internment policy displaced more than 100,000 Japanese Americans from their homes, 
primarily on the West Coast,
and placed them in camps scattered across remote areas of the Western US and Arkansas.
Following the war, 
efforts were made to disperse the Japanese American population more widely across the US,
often discouraging their return to prewar enclaves.
These policies highlight the need to examine the ways in which
forced relocation altered traditional patterns of settlement
and migration within this population.

To answer some of these questions, 
I examine how the proximity of a county to a WRA internment camp correlates with the migration rates of Japanese Americans 
relative to the total migration rates in the county.
My analysis draws on decennial census data from 1940 to 1990,
a period in which the effects of internment
can be observed in the migration choices of Japanese Americans.

In addition to proximity to internment camps,
this paper considers whether the effects of proximity to internment camps on migration rates 
differed for counties located within the "Evacuation Zone"
—the areas from which Japanese Americans were originally forced to leave.
Did counties that were themselves part of the evacuation orders 
exhibit different migration dynamics compared to those that were outside the zone?
By addressing these questions, I aim to provide insight into how the legacy of internment shaped not only immediate resettlement 
but also longer-term demographic shifts within the Japanese American community.

I aim to understand whether forced displacement left lasting scars on the geography of the Japanese American population,
and whether the government’s efforts to disperse internees resulted in a permanent reshaping of settlement patterns.
The findings from this analysis could help shed light on the broader implications of forced migration policies 
and provide lessons for understanding the long-term effects of similar interventions in other contexts.


\subsection{Motivating Literature}\label{motivation}

Several existing studies have tried to quantify the causal effect of Japanese internment on 
long-term earnings
\citep{chinLongRunLaborMarket2005}, 
childhood health
(\cite{saavedraEarlyChildhoodConditions2013}, \cite{grossmanIntergenerationalHealthEffects2025}),
and educational attainment 
\citep{saavedraSchoolQualityEducational2015}. 

On the other hand, other studies have identified surprising ways in which internees did better than their counterparts.
\cite{shoagCausalEffectPlace2016} examine the persistant pattern of internees' tendency to settle closer to their assigned camp,
and as a result find that differential exposure to more affluent areas lead to some internees experiencing more long-term upward mobility than others.
\cite{arellano-boverDisplacementDiversityMobility2022} finds that overall, interned Japanese men experienced 9-22\% higher incomes than similar Chinese and non-interned Japanese men.
He attributes this effect to a combination of information and skills transfers among those living in camps together,
migration to new labor markets, and career switching that would not have happened if internees had remained in their original homes on the West Coast.

Aside from the importance of understanding the losses to internees, Japanese
Internment can also serve as an interesting natural experiment through the
involuntary nature of the experiences of internees. 
A key question in previous economic literature is the extent to which 
geographic location influences economic outcomes, such as upward mobility or earnings.
It is difficult to provide an unbiased causal estimate because households usually
have a large degree of choice in where they live. Previous studies have used a
variety of empirical strategies to separate out this causal effect including
controlling for observable characteristics like education
\cite{glaeserCitiesSkills2001}, worker fixed-effects \citep{cardLocationLocationLocation2021}, or
through leveraging other exogenous placements of households like the Moving to
Opportunity Experiment (\cite{ludwigLongTermNeighborhoodEffects2013},
\cite{chettyEffectsExposureBetter2016}, \cite{bergmanCreatingMovesOpportunity2019}, etc). 